\section{Experimental design}
\paragraph{Compact code blocks.}
The main experiments use rotated $[[9,1,3]]$ and $[[25,1,5]]$ surface-code blocks and $[[72,12,6]]$ BB blocks. Four-block circuits monitor eight logical Bell outputs for the surface code and 96 for BB. An eight-block distance-five surface circuit monitors 16 outputs and exercises multiple repair seams. Every block begins in an ideal encoded Bell pair with noiseless references, followed by transversal CNOT matchings, noisy syndrome extraction, and ideal final closure~\cite{tour}. The decoder receives the complete syndrome record and partitions it by spatial fault ownership. Surface-code extraction uses explicit ancilla schedules with depolarizing gate noise, reset flips, and readout flips at probability $p$. The BB control uses data depolarization and readout-flip noise with ideal transversal CNOTs and ideal stabilizer-product measurements. Stim generates the exact undecomposed detector models~\cite{stim}. Appendix~\ref{app:protocol} gives all circuit configurations and implementation details.

\paragraph{Matched engine and policy comparisons.}
BF, BP+LSD0, BP+LSD2, and BP+OSD0 each run joint decoding, direct cutting, and radius-one repair on identical fault samples~\cite{bfdocs,lsddocs,osddocs}. They share 30 min-sum BP iterations, scaling 0.625, parallel message scheduling, and one native CPU thread. BF additionally tests halo, radius-two, and rank-screened variants. A radius-one band contains two owner blocks; every replacement conditions on its exterior and retains the complete touched-row detector constraints. In four-block circuits, radius two covers all nonzero detector columns; in eight-block circuits, its bands contain four owner blocks.

\paragraph{Paired inference and execution measurements.}
The complete original study contains 4,096 seeded circuit samples, 15 configurations, and 16 decoder stacks. The frozen protocol specifies nine contrasts per configuration. We use exact conditional McNemar tests~\cite{mcnemar} with Holm adjustment across all 135 contrasts~\cite{holm}. The plots show 95\% marginal Wilson intervals~\cite{wilson}. The full outcomes, effect intervals, seeds, and protocol history are retained in the original study. Engineering measurements use three fresh-process replicates on 64 common inputs, exclude startup from warm-call latency, and separately record startup and process-tree memory. Reused timing and functional-audit inputs do not increase the independent logical sample count.

\paragraph{Controlled backend and external-baseline study.}
The new study reuses stored circuits, detector matrices, and noise samples from the original repair-policy grid. It compares a rebuilding HUF reference with the persistent conditional backend inside the same ownership/cut/repair wrapper. Both use identical weighted boundary growth and terminal canonicalization. Tesseract~\cite{tesseract} and native BP+LSD2/OSD2 are evaluated on the same complete detector records. Appendix~\ref{app:study} specifies the pilot budgets, scaling circuits, timing scopes, and unavailable baselines.

\section{Compact-code results}\label{sec:results}
\subsection{Logical failure rates across inner decoders}
The primary benchmark uses four $[[25,1,5]]$ blocks at $p=0.006$. Radius-one boundary repair reduces BF failures from 65 to 45 in 256 shared shots, a 30.8\% relative reduction. The paired outcomes contain 21 corrected failures and one new failure, giving a Holm-adjusted $p=0.00147$. BP+LSD2 improves from 58 to 38 failures, a 34.5\% reduction with adjusted $p=0.00477$. All four cut-to-repair contrasts in Table~\ref{tab:primary} are significant after Holm adjustment across the 135 planned comparisons.

\begin{table*}[t]\centering\small
\caption{Logical failures before and after radius-one repair on 256 paired shots for four $[[25,1,5]]$ blocks at $p=0.006$. Paired gain is the absolute decrease in failure rate in percentage points (pp), with a conservative 95\% interval.}\label{tab:primary}
\begin{tabular}{lrrrrrr}\toprule
Engine&Cut&Repair&Fixed&New&Paired gain, pp [95\% CI]&Holm $p$\\\midrule
BF&65&45&21&1&7.8 [2.3, 12.9]&0.00147\\
LSD0&65&43&24&2&8.6 [2.6, 14.2]&0.00142\\
LSD2&58&38&22&2&7.8 [2.0, 13.3]&0.00477\\
OSD0&62&43&21&2&7.4 [1.7, 12.8]&0.00872\\\bottomrule
\end{tabular}\end{table*}
In the $d=5$, $p=0.003$ case, BF failures decrease from 55 to 39 on 1,024 paired shots. Figure~\ref{fig:failures} shows the surface-code and BB72 operating points with Wilson intervals.

\begin{figure*}[t]\centering\includegraphics[width=\textwidth]{figures/fig4.pdf}
\caption{BF direct cutting and radius-one repair on paired samples. The $d=3$ surface-code and BB72 cases use 512 and 128 shots per point, respectively. The $d=5$ cases use 1,024 shots at $p=0.003$ and 256 at $p=0.006$. Error bars show marginal 95\% Wilson intervals.}\label{fig:failures}\end{figure*}

\subsection{Extension to compact high-rate BB blocks}
The $[[72,12,6]]$ BB case extends boundary repair to blocks that each encode 12 logical qubits. At $p=0.015$, BP+LSD2 reduces the observed count from 44 to 30 failures in 128 shots, a 31.8\% reduction; BF reduces 44 to 32, a 27.3\% reduction. Each shot simultaneously monitors both Pauli sectors through 96 Bell-reference observables. These results describe the tested phenomenological BB noise model.

\subsection{Half-size bands with full logical reach}
The rank identity explains why a narrow owner band can address a rich set of circuit-level logical alternatives. In the four-block distance-five model, a radius-one band contains 5,299 of the full problem's 10,164 fault variables and retains all eight dimensions of monitored logical freedom. In BB72 it contains exactly half the fault variables and retains all 96 logical-change dimensions. The compact $d=3$ band accesses seven of eight dimensions. Table~\ref{tab:reach} makes this local-to-global relationship explicit.
\begin{table*}[t]\centering\small
\caption{Logical reach of a radius-one band in four-block circuits.}\label{tab:reach}
\begin{tabular}{lrrr}\toprule
Code block&Fault variables (local / full)&Share (\%)&$\rho$ (local / full)\\\midrule
Surface $[[9,1,3]]$&949 / 1,821&52.1&7 / 8\\
Surface $[[25,1,5]]$&5,299 / 10,164&52.1&8 / 8\\
BB $[[72,12,6]]$&3,456 / 6,912&50.0&96 / 96\\\bottomrule
\end{tabular}\end{table*}
Full logical reach in the distance-five and BB72 bands means that every monitored logical offset is attainable by a syndrome-consistent replacement with the exterior fixed. The inner decoder and acceptance rule determine which feasible replacements are made in practice. The nested-support measurements in Appendix~\ref{app:surface} resolve how that reach emerges as columns enter a band. Proposition~\ref{prop:exposed} additionally separates this full-band capacity from the support actually inserted by persistent HUF; the original rank table does not measure the latter.

\subsection{Parallel repair on eight compact surface-code blocks}
The eight-block $[[25,1,5]]$ circuit has three disjoint radius-one seam bands. Two worker processes execute those updates concurrently and return bitwise-identical corrections to the serial implementation. Mean warm-call latency decreases from 60.00 to 49.71 ms, corresponding to $1.21\times$ the serial-repair throughput. Direct cutting takes 35.44 ms. Parallel repair increases peak process-tree proportional set size (PSS) from 88.1 to 233.1 MiB relative to serial repair (Table~\ref{tab:parallel}). At distance three, parallel repair increases mean warm-call latency from 3.53 to 4.57 ms (Table~\ref{tab:engineering}). Warm-call timings pool three fresh-process replicates. PSS is sampled separately on the first eight inputs of each run. Appendix~\ref{app:engineering} gives the complete execution grid.

All six radius-one process-parallel trials reproduce serial corrections bit for bit. An additional 32-input functional audit executes overlapping radius-two bands in two colors and matches the serial reference with the same color order. The implementation therefore realizes both the disjoint-update composition theorem and a concrete conflict schedule.
\begin{table}[t]\centering\small
\caption{Warm-call performance and logical failures on 64 shared shots for eight $[[25,1,5]]$ blocks at $p=0.003$.}\label{tab:parallel}
\begin{tabular}{lrrrr}\toprule
Method&Failures&Mean ms&p95 ms&PSS MiB\\\midrule
Joint BF&9 / 64&68.78&149.46&74.6\\
Cut BF&9 / 64&35.44&77.20&82.3\\
Serial repair&3 / 64&60.00&137.54&88.1\\
Parallel repair&3 / 64&49.71&122.24&233.1\\\bottomrule
\end{tabular}\end{table}

\subsection{Controlled persistent-backend circuit pilot}\label{sec:newpilot}
Across 24 circuit configurations and 168 planned configuration/backend jobs, 136 jobs complete the full measurement protocol. The 37 matched rebuilding/persistent comparisons cover 2,560 paired shot executions, with zero correction-vector or recorded growth-history disagreements. Full details and non-completions are retained in the archive.
\begin{table*}[t]\centering\scriptsize
\caption{Completed controlled pilot runs. HUF/SW: complete Seamwise with prior-weighted rebuilding/persistent HUF. BPHUF/BPSW: matched BP-assisted counterparts. Tesseract, LSD2 and OSD2: native joint decoders. RSS is the separately sampled whole-process footprint, not backend-only storage.}\label{tab:newpilot}
\begin{tabular}{llrrrr}\toprule Configuration&Backend&Failures/shots&Mean ms&p95 ms&RSS MiB\\\midrule
bb144\_p0.005\_B4&BPHUF&0/64&134.74&209.08&136.4\\
bb144\_p0.005\_B4&BPSW&0/64&132.83&204.10&137.1\\
bb144\_p0.005\_B4&HUF&2/64&224.01&376.17&126.5\\
bb144\_p0.005\_B4&LSD2&0/64&13.46&21.00&74.7\\
bb144\_p0.005\_B4&OSD2&0/64&959.25&1151.50&175.5\\
bb144\_p0.005\_B4&SW&2/64&225.91&374.91&126.5\\
bb144\_p0.005\_B4&Tesseract&0/64&14.75&19.29&117.4\\
surface\_d3\_p0.003\_B4&BPHUF&2/64&1.99&6.91&75.6\\
surface\_d3\_p0.003\_B4&BPSW&2/64&2.05&6.99&75.4\\
surface\_d3\_p0.003\_B4&HUF&20/64&6.22&17.67&74.0\\
surface\_d3\_p0.003\_B4&LSD2&2/64&0.39&1.29&69.0\\
surface\_d3\_p0.003\_B4&OSD2&3/64&1.55&7.24&72.6\\
surface\_d3\_p0.003\_B4&SW&20/64&6.68&18.91&73.7\\
surface\_d3\_p0.003\_B4&Tesseract&2/64&2.15&8.39&97.1\\
surface\_d5\_p0.006\_B4&BPHUF&29/256&143.39&369.52&113.1\\
surface\_d5\_p0.006\_B4&BPSW&29/256&132.00&329.58&113.0\\
surface\_d5\_p0.006\_B4&LSD2&47/256&15.05&47.81&73.2\\
surface\_d5\_p0.006\_B4&Tesseract&10/256&169.88&432.11&113.7\\
\bottomrule\end{tabular}\end{table*}
\begin{figure*}[t]\centering\includegraphics[width=\textwidth]{study_scaling.pdf}\caption{Measured scaling pilot at fixed noise p=.003. Block sweep fixes d=3 and T=3; distance sweep fixes four blocks and T=3; depth sweep fixes d=3 and four blocks. Error bars are marginal Wilson intervals. Only completed, fixed-sample runs are plotted; fixed shallow extraction depth is not a threshold experiment.}\label{fig:newscale}\end{figure*}
\begin{figure*}[t]\centering\includegraphics[width=\textwidth]{study_clusters.pdf}\caption{Peak-cluster distributions over recorded completed solver calls. Each panel uses configurations completed by both members of that oracle pair. BP early exits have zero peak; infeasible rejected-call histories are unavailable. The two panels need not have the same configuration cohort.}\label{fig:newclusters}\end{figure*}
The controlled study does not establish a general latency advantage of the current persistent implementation. The exact-oracle guarantee concerns reusable algebra; full-call time additionally includes boundary scans, state creation, canonicalization and wrapper costs. Native external baselines test the remaining implementation and accuracy gap. Small pilot sample sizes, censored non-completions, and untuned resource settings limit interpretation to the measured implementation and workloads.
\subsection{Independent-process timing replication}
Three fresh-process replicates follow the controlled pilot at the primary surface-code point and the high-noise BB144 point. Each uses the same first 64 inputs with randomized method order. These are additional timing executions, not new independent logical samples. All repeated predictions match the corresponding stored single-pass predictions.
\begin{table*}[t]\centering\scriptsize
\caption{Three independently initialized timing replicates on 64 reused inputs per operating point. RSS is the maximum separately sampled whole-process footprint across replicates.}\label{tab:replication}
\begin{tabular}{llrrrrr}\toprule Point&Backend&Run 1 mean ms&Run 2&Run 3&Pooled p95 ms&Max RSS MiB\\\midrule
BB144&BPHUF&347.53&347.53&348.37&601.85&137.7\\
BB144&BPSW&285.45&288.64&283.37&454.90&137.7\\
BB144&Tesseract&55.98&55.41&59.46&69.05&117.1\\
Surface d5&BPHUF&132.92&130.67&133.22&343.95&112.4\\
Surface d5&BPSW&122.26&119.00&123.61&328.14&111.8\\
Surface d5&Tesseract&117.66&113.71&115.00&246.90&106.8\\
\bottomrule\end{tabular}\end{table*}
The median paired reduction in mean latency is 8.0\% at the surface-code point and 17.9\% at the BB144 point. The corresponding whole-process RSS measurements show little change between oracle modes. Native Tesseract remains a stronger accuracy comparator on these samples; the timing table also exposes dependence on which fixed sample subset is measured.
A separate diagnostic on eight fixed high-noise BB144 inputs records 81,639 rebuilding basis-insertion calls and 3,100 persistent calls, a 26.3-fold reduction in this repeated-column-processing counter. Both diagnostic outputs are identical. The instrumented times are excluded from benchmark latency. Profiling identifies repeated component creation, DSU root queries and eligibility scans as substantial remaining costs; those paths are preserved by the current oracle replacement.



\section{Conclusion}
Seamwise combines fixed-exterior boundary repair, logical-reach analysis, merge-persistent conditional solving, and exact scheduling of disjoint updates. Accepted repairs preserve the full detector record and decrease the prior configuration cost. The persistent state eliminates repeated reconstruction of inherited detector relations, while the exposed-reach certificate characterizes the logical directions represented by the inserted fault support.

The original compact-circuit policy study demonstrates 30.6--34.5\% fewer logical failures than direct cutting in its primary comparison. The new controlled study tests the implementation cost of the persistent backend and compares full-record native baselines on identical noise samples. Exact output preservation is an algorithmic guarantee under canonicalization; a practical speed advantage and competitive logical accuracy require measurement. These results characterize the measured Python reference and circuit sizes. Native optimized implementations, common-budget accuracy--latency frontiers, and complete 2026 comparator coverage remain to be evaluated.
